\documentclass[11pt]{article}
\usepackage{amsmath,amssymb,amsthm,mathtools}
\usepackage[margin=1in]{geometry}
\theoremstyle{plain}
\newtheorem{theorem}{Theorem}
\newtheorem{lemma}{Lemma}
\theoremstyle{remark}
\newtheorem{remark}{Remark}
\newcommand{\E}{\mathbb{E}}\newcommand{\PP}{\mathbb{P}}
\newcommand{\To}{\Rightarrow}

\title{Formalization brief: the abstract two--phase mixture CLT\\
(the array core of \texttt{prop:twophase})}
\date{5 July 2026}

\begin{document}
\maketitle

\noindent\textbf{Goal.} A new self-contained file (suggest
\texttt{TypeDDecouplingTwoPhase.lean}) proving the abstract two--phase central
limit theorem whose model instantiation is
\texttt{TypeDDecoupling.prop\_twophase} in
\texttt{TypeDDecouplingCrossover.lean}. Build on
\texttt{TypeDDecouplingMartingaleCLT.lean} (attached): reuse
\texttt{core\_charFun\_tendsto} / \texttt{joint\_charFun\_tendsto}; do not
duplicate their proofs. Do NOT touch \texttt{prop\_twophase}'s statement or its
\texttt{sorry}; a docstring cross-reference is the only permitted change to the
Crossover file. Whole project must build; standard axioms only.

\section*{Setting}

For each $n$: a filtration $(\mathcal F_{n,j})_j$, a pair of square-integrable
arrays $(X_{n,j},Y_{n,j})_{j=1}^{k_n}$ adapted with the martingale difference
property $\E[X_{n,j}\mid\mathcal F_{n,j-1}]=\E[Y_{n,j}\mid\mathcal F_{n,j-1}]=0$,
jumps deterministically bounded: $|X_{n,j}|,|Y_{n,j}|\le b_n\to0$. A
\underline{change point} $m_n\in\{0,\dots,k_n\}$.

\underline{Phase 1 (locked):} for $j\le m_n$, $X_{n,j}=Y_{n,j}$
(the two components receive identical increments --- the paper's bound pair
moving as one particle before the split).

\underline{Phase 2 (decoupled):} for $j>m_n$, the brackets are diagonal in the
limit.

\section*{Tier 1 (primary): fixed change-point fraction}

\begin{theorem}[two--phase CLT at deterministic $u$]\label{thm:twophase-u}
Fix $u\in[0,1]$ and suppose $m_n/k_n\to u$ together with (all limits a.e.\ or
in the a.e.-bounded form used in \textup{\texttt{MartingaleCLT}}):
\[
  \sum_{j\le m_n}X_{n,j}^2\to u,\qquad
  \sum_{j>m_n}X_{n,j}^2\to 1-u,\qquad
  \sum_{j>m_n}Y_{n,j}^2\to 1-u,\qquad
  \sum_{j>m_n}X_{n,j}Y_{n,j}\to 0 .
\]
Then, with $S_n=\sum_jX_{n,j}$, $R_n=\sum_jY_{n,j}$, for all $a,c\in\R$,
\[
  \E\big[e^{\mathrm i(aS_n+cR_n)}\big]\ \longrightarrow\
  \exp\Big(-\tfrac12\,\psi_u(a,c)\Big),\qquad
  \psi_u(a,c):=(a+c)^2u+(a^2+c^2)(1-u),
\]
the characteristic function of
$\big(\sqrt u\,\xi_0+\sqrt{1-u}\,\xi_1,\ \sqrt u\,\xi_0+\sqrt{1-u}\,\xi_2\big)$
with $\xi_0,\xi_1,\xi_2$ i.i.d.\ $N(0,1)$ --- the standard bivariate normal of
correlation $u$.
\end{theorem}

\noindent\textbf{Route (short).} Cram\'er--Wold through the existing Tier-1
theorem: for fixed $(a,c)$ the combined array $Z_{n,j}:=aX_{n,j}+cY_{n,j}$ is a
martingale difference array with jumps $\le(|a|+|c|)b_n\to0$ and bracket
\[
  \sum_jZ_{n,j}^2
  =(a+c)^2\!\sum_{j\le m_n}\!X_{n,j}^2
  +a^2\!\sum_{j>m_n}\!X_{n,j}^2+c^2\!\sum_{j>m_n}\!Y_{n,j}^2
  +2ac\!\sum_{j>m_n}\!X_{n,j}Y_{n,j}
  \ \to\ \psi_u(a,c),
\]
using $X=Y$ in phase 1; then apply
\texttt{core\_charFun\_tendsto} at $\sigma^2=\psi_u(a,c)$ and evaluate its
conclusion at $1$ (or at $(a,c)$ via the same device as
\texttt{joint\_charFun\_tendsto}). Please also record the two edge cases
$u=0$ (independent pair) and $u=1$ (identical pair) as instances.

\section*{Tier 2: the mixture over a random change point}

\begin{lemma}[mixture lemma]\label{lem:mix}
Let $Z_n$ be $\R^2$-valued random vectors, $\mathcal G_n$ sub-$\sigma$-algebras,
$U_n$ $\mathcal G_n$-measurable with $U_n\to U$ a.e., $U$ taking values in
$[0,1]$. Suppose for each fixed $(a,c)$
\[
  \E\big[e^{\mathrm i(aZ_n^{(1)}+cZ_n^{(2)})}\,\big|\,\mathcal G_n\big]
  \;-\;\exp\big(-\tfrac12\psi_{U_n}(a,c)\big)\ \longrightarrow\ 0
  \quad\text{a.e.\ (or in } L^1).
\]
Then
$\E[e^{\mathrm i(aZ_n^{(1)}+cZ_n^{(2)})}]\to
\E\big[\exp(-\tfrac12\psi_U(a,c))\big]$: the limit is the $U$-mixture of
correlation-$U$ bivariate normals --- the law of
$\big(\sqrt U\xi_0+\sqrt{1-U}\xi_1,\ \sqrt U\xi_0+\sqrt{1-U}\xi_2\big)$ with
$U$ independent of the $\xi$'s.
\end{lemma}

\noindent Route: tower property + dominated convergence (everything is bounded
by $1$; $u\mapsto\psi_u(a,c)$ is continuous). This is deliberately elementary;
the force of the hypothesis (the conditional Tier-1 applicability) remains with
the caller. Also record the specialisation of the limit functional at
$U\sim\min(\mathrm{Exp}(\lambda),1)$:
$\E[U]=(1-e^{-\lambda})/\lambda$ (so the mixture's covariance is
$\E[U]$, matching \texttt{thm\_closed}'s value at $\lambda=4c$) --- if
covariance extraction from the charFun is awkward, state the $\E[U]$
computation as a standalone lemma about $\min(\mathrm{Exp}(\lambda),1)$ and
connect by docstring.

\section*{Tier 3 (stretch): assembled statement}

Combine Tiers 1--2: a two--phase array with a $\mathcal G_n$-measurable random
change point $M_n$, $M_n/k_n\to U$ a.e., whose conditional brackets satisfy the
Tier-1 hypotheses conditionally on $\mathcal G_n$ (as a named hypothesis if the
conditional formalization of \texttt{core\_charFun\_tendsto} would require
re-proving it conditionally --- do NOT re-prove; an unconditional-on-events or
hypothesis-based route is sanctioned), concluding the mixture charFun limit.
Then a docstring in the Crossover file recording: the probabilistic content of
\texttt{prop\_twophase} (two-phase decomposition at the split time, mixture
over $\tau/T\to U\sim\min(\mathrm{Exp}(4c),1)$) is now the assembled theorem;
what remains in the leaf is the model bookkeeping (identifying the compensated
jump martingales' discretized increments with the array and the bracket
occupation estimates of \texttt{lem:occ}/\texttt{lem:rebind}).

\begin{remark}[interface and fallbacks]
(1) Reuse \texttt{MartingaleCLT} lemmas; if a hypothesis there is stated
slightly differently than convenient (e.g.\ the bracket convergence form),
add adapter lemmas in the new file rather than modifying the old one.
(2) Sanctioned fallbacks: conditional-expectation hypotheses may be taken in
$L^1$ or a.e.\ form, whichever is smoother; the mixture lemma may assume $U_n,U$
take values in a compact interval (they do); Tier 3 may package the conditional
Tier-1 hypothesis as one named hypothesis.
(3) Degenerate cases: $\psi_u$ is a positive semidefinite quadratic form which
degenerates at $(a,-a)$, $u=1$; \texttt{core\_charFun\_tendsto} at
$\sigma^2=0$ (if it requires $\sigma^2>0$, handle $\psi=0$ directly: the sum
converges to $0$ in $L^2$, so the charFun tends to $1$).
\end{remark}

\end{document}
